\title{\textbf{Are Prediction Markets Accurate? Resolving Information Aggregation Through a Lifecycle Lens}}
\author{Jonathan Becker\thanks{Email: \texttt{jonathan@jbecker.dev}} \\ \vspace{0.3cm} \textit{Independent Researcher}}
\date{\today}

\maketitle

\begin{abstract}

Prediction markets exhibit strong calibration on average, but accuracy varies dramatically over a market's lifetime. A contract trading at 50 cents wins approximately 50\% of the time, confirming efficient information incorporation. However, markets at inception show 5.4$\times$ worse calibration (mean absolute error 22.56\%) than mature markets (4.18\% MAE). This paper introduces a lifecycle model grounding these patterns in Ottaviani \& Sørensen's (2007, 2010) theory of wealth-constrained information aggregation.

We analyze 7.3 million finalized Kalshi markets covering 72 million trades (October 2021--November 2025)---the largest empirical sample in prediction market literature---to document three robust associations: (1) a liquidity threshold at $\sim$200 trades marking the transition from thin to liquid regimes, (2) a 4.2$\times$ volume association with accuracy, and (3) market-age calibration decay with confounding from information arrival. Robustness checks across eight specifications, placebo tests, and Heckman selection-bias analysis strengthen causal arguments. These findings have immediate implications for market designers (seeding capital to reach $\sim$200 trades is critical) and forecast users (markets with $<$100 trades should be treated as exploratory).

\noindent\textbf{Keywords:} prediction markets, calibration, information aggregation, liquidity, lifecycle model, Ottaviani-Sørensen, Kalshi, wealth constraints

\end{abstract}
